\subsection{定向族；测度的支集}

命题 5. --- a) 设 H 是由在 T 的每个紧子集上下半连续的 \(\geqslant 0\) 的函数组成的递增定向集。则

\[
\mu^ {\bullet} \left(\sup _ {h \in \mathrm{H}} h\right) = \sup _ {h \in \mathrm{H}} \mu^ {\bullet} (h).\tag{5}
\]

\begin{enumerate}
\def\labelenumi{\alph{enumi})}
\setcounter{enumi}{1}
\tightlist
\item
  设 H 是由在 T 的每个紧子集上上半连续的 \(\geqslant0\) 的函数组成的递减定向集。若 H 中存在函数 \(h_{0}\) 使 \(\mu^{\bullet}(h_{0})<+\infty\) ，则
\end{enumerate}

\[
\mu^ {\bullet} \left(\inf _ {h \in \mathrm{H}} h\right) = \inf _ {h \in \mathrm{H}} \mu^ {\bullet} (h).\tag{6}
\]

对每个紧集 \(\mathbf{K} \subset \mathbf{T}\) ，在情形 \(a\) ) 我们有

\[
\mu^ {\bullet} \Big (\sup _ {h \in \mathrm{H}} h \varphi_ {\mathrm{K}} \Big) = \mu_ {\mathrm{K}} ^ {\bullet} \Big (\sup _ {h \in \mathrm{H}} h _ {\mathrm{K}} \Big) = \sup _ {h \in \mathrm{H}} \mu_ {\mathrm{K}} ^ {\bullet} (h _ {\mathrm{K}}) = \sup _ {h \in \mathrm{H}} \mu^ {\bullet} (h \varphi_ {\mathrm{K}}),
\]

而在情形 \(b\)

\[
\mu^ {\bullet} \Big (\inf _ {h \in \mathrm{H}} h \varphi_ {\mathrm{K}} \Big) = \mu_ {\mathrm{K}} ^ {\bullet} \Big (\inf _ {h \in \mathrm{H}} h _ {\mathrm{K}} \Big) = \inf _ {h \in \mathrm{H}} \mu_ {\mathrm{K}} ^ {\bullet} (h _ {\mathrm{K}}) = \inf _ {h \in \mathrm{H}} \mu^ {\bullet} (h \varphi_ {\mathrm{K}}),
\]

（由 No.~2 的 Prop. 2 以及 Ch. V, §1, No.~2 的 Prop. 8）。情形 a) 通过关于 K 取上确界 (No.~2, Prop. 2) 立即得出。为处理情形 b)，用 \(\varepsilon\) 表示一个数 \(>0\) ，并选取紧集 K 使 \(\mu^{\bullet}(h_0\varphi_{\mathbf{K}}) \geqslant \mu^{\bullet}(h_0) - \varepsilon\) 。于是有 (No.~5, Prop. 4) \(\mu^{\bullet}(h_0\varphi_{\mathbf{C}_{\mathbf{K}}}) \leqslant \varepsilon\) ；故对每个函数 \(h \in \mathbf{H}\) ，只要它 \(\leqslant h_0\) ，就有 \(\mu^{\bullet}(h\varphi_{\mathbf{C}_{\mathbf{K}}}) \leqslant \varepsilon\) ，最后由 No.~5 的 Prop. 4 得 \(\mu^{\bullet}(h\varphi_{\mathbf{K}}) \geqslant \mu^{\bullet}(h) - \varepsilon\) 。于是

\[
\mu^ {\bullet} \left(\inf _ {h \in \mathrm{H}} h\right) \geqslant \mu^ {\bullet} \left(\inf _ {h \in \mathrm{H}} h \varphi_ {\mathrm{K}}\right) = \inf _ {h \in \mathrm{H}, h \leqslant h _ {0}} \mu^ {\bullet} (h \varphi_ {\mathrm{K}}) \geqslant \inf _ {h \in \mathrm{H}, h \leqslant h _ {0}} \mu^ {\bullet} (h) - \varepsilon .
\]

因此 (6) 的左端 \(\geqslant\) 右端；反向不等式是明显的，命题得证。

推论. --- a) 设 \((\mathrm{U}_{\alpha})_{\alpha \in \mathrm{I}}\) 是 T 的开子集的一个递增定向族，其并为 U。则 \(\mu^{\bullet}(\mathrm{U}) = \sup_{\alpha \in \mathrm{I}} \mu^{\bullet}(\mathrm{U}_{\alpha})\) 。

\begin{enumerate}
\def\labelenumi{\alph{enumi})}
\setcounter{enumi}{1}
\tightlist
\item
  设 \((\mathbf{F}_{\alpha})_{\alpha \in \mathbf{I}}\) 是 \(\mathbf{T}\) 的闭子集的一个递减定向族，其交为 \(\mathbf{F}\) 。若存在 \(\alpha \in \mathbf{I}\) 使 \(\mu^{\bullet}(\mathbf{F}_{\alpha})\) 有限，则 \(\mu^{\bullet}(\mathbf{F}) = \inf_{\alpha \in \mathbf{I}} \mu^{\bullet}(\mathbf{F}_{\alpha})\) 。
\end{enumerate}

由前面的推论，存在一个最大的局部可忽略开集；这证明了下述定义的合理性：

定义 8. --- T 上测度 \(\mu\) 的支集定义为 T 中最大的局部 \(\mu\) -可忽略开集的余集。

\(\mu\) 的支集记作 \(\operatorname{Supp}(\mu)\) 。

注记. --- 1) 若 \(\mu\) 是一个复测度，则 \(\mu\) 的支集定义为正测度 \(|\mu|\) 的支集；它仍是最大的局部 \(\mu\) -可忽略开集的余集。

\begin{enumerate}
\def\labelenumi{\arabic{enumi})}
\setcounter{enumi}{1}
\tightlist
\item
  我们来证明 No.~3 的例 2 中引入的测度是 T 上具有紧支集的测度。设 \(\mu\) 是 T 上的一个正测度，其支集是紧集 K，\(\nu\) 是由 \(\mu_{\mathbf{K}}\) 定义的测度（No.~3 的意义下）。设 \(f \in \mathcal{F}_{+}(\mathbf{T})\) ；则
\end{enumerate}

\[
\nu^ {\bullet} (f) = \mu_ {K} ^ {\bullet} (f _ {K}) \quad (\text { No.   3,   formula   (3)) }.
\]

由于负担 \(\mu^{\bullet}\) 集中于 K，我们还有

\[
\mu^ {\bullet} (f) = \mu^ {\bullet} (f \varphi_ {K}) = \mu^ {\bullet} ((f _ {K}) ^ {0}) = \mu_ {K} ^ {\bullet} (f _ {K}),
\]

故 \(\mu^{\bullet} = \nu^{\bullet}\) ，最后 \(\mu = \nu\) 。反之，若 K 是 T 中的紧集，\(\lambda\) 是 K 上的一个测度，而 \(\mu\) 是由 \(\lambda\) 定义的 T 上的测度，则 \(\mu^{\bullet}(\mathbb{C}K) = 0\) (No.~3, 公式 (3))；于是 \(\mu\) 的支集含于 K，从而是紧的。

